\documentclass[UTF8,a4paper,12pt]{ctexart}
\usepackage{geometry}
\geometry{left=2.4cm,right=2.4cm,top=2.5cm,bottom=2.5cm}
\usepackage{amsmath,amssymb,bm}
\allowdisplaybreaks[2]
\usepackage{xcolor}
\usepackage{fancyhdr}
\pagestyle{fancy}
\fancyhf{}
\fancyfoot[C]{\small\thepage}
\renewcommand{\headrulewidth}{0pt}
\renewcommand{\footrulewidth}{0pt}
\usepackage[hidelinks]{hyperref}
\hypersetup{
  pdftitle  = {CS2601 Homework 1},
  pdfauthor = {Li Jinshuo},
  pdfsubject = {Linear and Convex Optimization},
}

% Default fonts are kept on purpose (no newtxtext / newtxmath).

\ctexset{
  section = {format = {\centering\Large\bfseries}},
}

\newcommand{\R}{\mathbb{R}}
\newcommand{\intset}{\operatorname{int}}

\begin{document}

\begin{center}
  {\LARGE\bfseries CS2601 Linear and Convex Optimization}\\[0.5em]
  {\large Homework 1}\\[0.9em]
  {\normalsize Li Jinshuo, IEEE Honor Class}\\[0.35em]
  {\small Due: 2026.10.20}
\end{center}
\vspace{-0.2em}
{\centering\textcolor{black!30}{\rule{0.9\textwidth}{0.8pt}}\par}
\vspace{0.8em}

%========================================================
\section*{Problem 1}

\subsection*{(a)}
For
\[
f(x_1,x_2)=3x_1^2-2x_1x_2+4x_2^2+x_1-6x_2,
\]
the partial derivatives are
\[
\frac{\partial f}{\partial x_1}=6x_1-2x_2+1,
\qquad
\frac{\partial f}{\partial x_2}=-2x_1+8x_2-6 .
\]
Hence
\[
\boxed{
\nabla f(x)=
\begin{bmatrix}
6x_1-2x_2+1\\[2pt]
-2x_1+8x_2-6
\end{bmatrix}
}
\]
and, differentiating once more,
\[
\boxed{
\nabla^2 f(x)=
\begin{bmatrix}
6 & -2\\[2pt]
-2 & 8
\end{bmatrix}
}
\]
which is the constant matrix
$\begin{bmatrix}6&-2\\-2&8\end{bmatrix}$ (it does not depend on $x$).

\subsection*{(b)}
Expand the squared norm:
\[
f(w)=\lVert Xw-y\rVert_2^2+\lambda w^T\Lambda w
=w^TX^TXw-2y^TXw+y^Ty+\lambda w^T\Lambda w .
\]
Using $\nabla_w(w^TAw)=2Aw$ for symmetric $A$ and $\nabla_w(b^Tw)=b$,
\[
\boxed{
\nabla f(w)=2X^TXw-2X^Ty+2\lambda\Lambda w
=2\left(X^TX+\lambda\Lambda\right)w-2X^Ty
}
\]
and
\[
\boxed{
\nabla^2 f(w)=2\left(X^TX+\lambda\Lambda\right)
}
\]
Note both the gradient and the Hessian hold for every $w$, and the Hessian is
constant in $w$.

%========================================================
\section*{Problem 2}

\textbf{Yes, it always has a global minimum, and the minimum is unique.}

The Hessian obtained in Problem 1(b) is
\[
\nabla^2 f(w)=2\left(X^TX+\lambda\Lambda\right).
\]
$X^TX$ is positive semidefinite because $v^TX^TXv=\lVert Xv\rVert_2^2\ge0$ for
all $v$. Since $\Lambda$ is positive definite and $\lambda>0$, the matrix
$\lambda\Lambda$ is positive definite. Therefore for every $v\neq0$,
\[
v^T\left(X^TX+\lambda\Lambda\right)v
=\underbrace{\lVert Xv\rVert_2^2}_{\ge0}
+\underbrace{\lambda\,v^T\Lambda v}_{>0}
>0,
\]
so $X^TX+\lambda\Lambda$ is positive definite and in particular invertible.
Thus $f$ is strictly convex and its Hessian is uniformly positive definite;
equivalently, $f(w)\to+\infty$ as $\lVert w\rVert\to\infty$ (it is coercive), so
a global minimum exists and is unique.

Setting the gradient to zero,
\[
\nabla f(w)=2\left(X^TX+\lambda\Lambda\right)w-2X^Ty=0,
\]
we obtain the closed-form solution
\[
\boxed{
w^\star=\left(X^TX+\lambda\Lambda\right)^{-1}X^Ty
}
\]

%========================================================
\section*{Problem 3}

Write $z_i=y_ix_i^Tw$ and $g(z)=\log(1+e^{-z})$, so that
$f(w)=\sum_{i=1}^m g(z_i)$. The derivatives of $g$ are
\[
g'(z)=-\frac{e^{-z}}{1+e^{-z}}=-\frac{1}{1+e^{z}},
\qquad
g''(z)=\frac{e^{-z}}{\left(1+e^{-z}\right)^2}
=\sigma(z)\,\sigma(-z)>0,
\]
where $\sigma(z)=\dfrac{1}{1+e^{-z}}$ is the sigmoid function (note
$g'(z)=-\sigma(-z)$ and $\dfrac{e^{-z}}{1+e^{-z}}=\sigma(-z)$).

Since $\nabla_w z_i=y_ix_i$, the chain rule gives
\[
\boxed{
\nabla f(w)=\sum_{i=1}^m g'(z_i)\,y_ix_i
=-\sum_{i=1}^m \frac{y_ix_i}{1+e^{\,y_ix_i^Tw}}
=-\sum_{i=1}^m \sigma\!\left(-y_ix_i^Tw\right)y_ix_i
}
\]
Differentiating once more, using $y_i^2=1$ and the fact that
$\nabla_w z_i=y_ix_i$,
\[
\boxed{
\nabla^2 f(w)=\sum_{i=1}^m g''(z_i)\,y_i^2\,x_ix_i^T
=\sum_{i=1}^m \frac{e^{-y_ix_i^Tw}}
{\left(1+e^{-y_ix_i^Tw}\right)^2}\,x_ix_i^T
}
\]
Each term has a non-negative scalar coefficient multiplying the outer product
$x_ix_i^T$, hence every term is positive semidefinite and so is the sum:
$\nabla^2 f(w)\succeq0$. Therefore $f$ is convex (this is the standard
convexity of the logistic-regression loss).

%========================================================
\section*{Problem 4}

Each matrix below is symmetric. For a symmetric matrix the signs of its
principal minors settle its definiteness: it is positive definite exactly when
every leading principal minor is positive, and positive semidefinite exactly
when every principal minor is non-negative. Whenever the minors (or the
quadratic form) take both signs, the matrix is indefinite. We carry out the
computation for each case directly below.

\subsection*{Matrix $A$}
\[
A=\begin{bmatrix}
4 & 1 & 0\\
1 & 3 & -1\\
0 & -1 & 2
\end{bmatrix}.
\]
Leading principal minors:
\[
D_1=4>0,\qquad
D_2=\det\begin{bmatrix}4&1\\1&3\end{bmatrix}=12-1=11>0,
\]
\[
D_3=\det A
=4\bigl(3\cdot2-(-1)(-1)\bigr)-1\bigl(1\cdot2-(-1)\cdot0\bigr)+0
=4(5)-2=18>0 .
\]
All leading principal minors are positive, hence
\[
\boxed{A\ \text{is positive definite.}}
\]

\subsection*{Matrix $B$}
\[
B=\begin{bmatrix}
1 & 2 & 0\\
2 & 1 & 1\\
0 & 1 & 1
\end{bmatrix}.
\]
The leading principal minors are
\[
D_1=1>0,\qquad
D_2=\det\begin{bmatrix}1&2\\2&1\end{bmatrix}=1-4=-3<0,
\]
and
\[
\det B=1(1\cdot1-1\cdot1)-2(2\cdot1-1\cdot0)+0=-4\neq0 .
\]
$D_2<0$ rules out positive definiteness, and $D_1>0$ rules out negative
definiteness (which would require $D_1<0$). Since $\det B=-4\neq0$ the matrix is
not semidefinite either. Concretely, the quadratic form changes sign:
\[
x^TBx=x_1^2+x_2^2+x_3^2+4x_1x_2+2x_2x_3,
\]
with $x=(1,0,0)\Rightarrow x^TBx=1>0$ and
$x=(1,-1,0)\Rightarrow x^TBx=2-4=-2<0$. Therefore
\[
\boxed{B\ \text{is indefinite.}}
\]

\subsection*{Matrix $C$}
\[
C=\begin{bmatrix}
1 & 1 & 0\\
1 & 1 & 0\\
0 & 0 & 2
\end{bmatrix}.
\]
Its quadratic form factors as
\[
x^TCx=x_1^2+2x_1x_2+x_2^2+2x_3^2=(x_1+x_2)^2+2x_3^2\ge0
\]
for all $x$, so $C$ is positive semidefinite. It is \emph{not} positive definite
because the non-zero vector $x=(1,-1,0)$ gives $x^TCx=0$. Equivalently the
eigenvalues are $2,2,0$. Hence
\[
\boxed{C\ \text{is positive semidefinite (but not positive definite).}}
\]

%========================================================
\section*{Problem 5}

The matrix
\[
A=\begin{bmatrix}
3 & -1 & \alpha\\
-1 & 2 & -1\\
\alpha & -1 & 3
\end{bmatrix}
\]
is symmetric, so it is positive semidefinite iff all its principal minors are
non-negative.

The $1\times1$ principal minors are the diagonal entries $3,2,3$, all positive.

The $2\times2$ principal minors are
\[
\det\begin{bmatrix}3&-1\\-1&2\end{bmatrix}=6-1=5>0,
\qquad
\det\begin{bmatrix}2&-1\\-1&3\end{bmatrix}=6-1=5>0,
\]
\[
\det\begin{bmatrix}3&\alpha\\\alpha&3\end{bmatrix}=9-\alpha^2\ \ge0
\iff -3\le\alpha\le3 .
\]
The $3\times3$ principal minor is $\det A$. Expanding along the first row,
\[
\det A
=3(6-1)-(-1)(-3+\alpha)+\alpha(1-2\alpha)
=15-3+\alpha+\alpha-2\alpha^2
=12+2\alpha-2\alpha^2 .
\]
Requiring $\det A\ge0$:
\[
12+2\alpha-2\alpha^2\ge0
\iff \alpha^2-\alpha-6\le0
\iff (\alpha-3)(\alpha+2)\le0
\iff -2\le\alpha\le3 .
\]
Combining this with $-3\le\alpha\le3$ gives the final condition
\[
\boxed{-2\le\alpha\le3 .}
\]

%========================================================
\section*{Problem 6}

Let $C_1,C_2\subseteq\R^n$ be convex. Take any two points
$z_1,z_2\in C_1+C_2$. By definition of the Minkowski sum there exist
$x_1,x_2\in C_1$ and $y_1,y_2\in C_2$ with
\[
z_1=x_1+y_1,\qquad z_2=x_2+y_2 .
\]
Let $\theta\in[0,1]$ and set $\bar\theta=1-\theta$. Then
\[
\theta z_1+\bar\theta z_2
=\theta(x_1+y_1)+\bar\theta(x_2+y_2)
=\underbrace{\bigl(\theta x_1+\bar\theta x_2\bigr)}_{\in\,C_1}
+\underbrace{\bigl(\theta y_1+\bar\theta y_2\bigr)}_{\in\,C_2},
\]
where $\theta x_1+\bar\theta x_2\in C_1$ and
$\theta y_1+\bar\theta y_2\in C_2$ by convexity of $C_1$ and $C_2$. Their sum
therefore belongs to $C_1+C_2$. Since $z_1,z_2\in C_1+C_2$ and
$\theta\in[0,1]$ were arbitrary, the Minkowski sum is convex.

%========================================================
\section*{Problem 7}

Let $C\subseteq\R^n$ be convex.

\subsection*{(a) The interior $\intset C$ is convex.}

Take $x_1,x_2\in\intset C$ and $\theta\in[0,1]$, and put
$x=\theta x_1+\bar\theta x_2$ with $\bar\theta=1-\theta$. Since $x_1,x_2$ are
interior points, there exist $r_1,r_2>0$ such that
$B(x_1,r_1)\subseteq C$ and $B(x_2,r_2)\subseteq C$. Set
$r=\min(r_1,r_2)>0$.

We claim $B(x,r)\subseteq C$. Let $z\in B(x,r)$, i.e.\ $\lVert z-x\rVert<r$, and
write $v=z-x$. Then
\[
z=x+v=\theta(x_1+v)+\bar\theta(x_2+v).
\]
Now
\[
\lVert(x_1+v)-x_1\rVert=\lVert v\rVert<r\le r_1
\ \Longrightarrow\ x_1+v\in B(x_1,r_1)\subseteq C,
\]
and similarly $x_2+v\in B(x_2,r_2)\subseteq C$. Since $C$ is convex, the convex
combination $z=\theta(x_1+v)+\bar\theta(x_2+v)$ also lies in $C$. Hence
$B(x,r)\subseteq C$, so $x\in\intset C$. Therefore $\intset C$ is convex.

\subsection*{(b) The closure $\bar C$ is convex.}

Take $x_1,x_2\in\bar C$ and $\theta\in[0,1]$, and set
$x=\theta x_1+\bar\theta x_2$ with $\bar\theta=1-\theta$. By the sequential
characterization of the closure, there exist sequences
$(x_1^{(k)})_{k\ge1}\subseteq C$ and $(x_2^{(k)})_{k\ge1}\subseteq C$ with
\[
x_1^{(k)}\to x_1,\qquad x_2^{(k)}\to x_2 .
\]
For every $k$, convexity of $C$ gives
$\theta x_1^{(k)}+\bar\theta x_2^{(k)}\in C$. Passing to the limit and using the
continuity of the linear combination,
\[
\theta x_1^{(k)}+\bar\theta x_2^{(k)}
\ \longrightarrow\
\theta x_1+\bar\theta x_2=x .
\]
Thus $x$ is the limit of a sequence in $C$, so $x\in\bar C$. Since $x_1,x_2\in
\bar C$ and $\theta\in[0,1]$ were arbitrary, $\bar C$ is convex.

\end{document}
